% !TeX spellcheck = en_US
\documentclass[french]{yLectureNote}

\title{Électromagnétisme 1 PS}
\subtitle{Physique}
\author{Paulhenry Saux}
\date{\today}
\yLanguage{Français}
%lionels.calmels@cemes.fr
\professor{F.Pettinari}
\usepackage{graphicx}%----pour mettre des images
\usepackage[utf8]{inputenc}%---encodage
\usepackage{geometry}%---pour modifier les tailles et mettre a4paper
%\usepackage{awesomebox}%---pour les boites d'exercices, de pbq et de croquis ---d\'esactiv\'e pour les TP de PC
\usepackage{tikz}%---pour deiffner + d\'ependance de chemfig
% \usepackage{tabularx}%---pour dimensionner automatiquement les tableaux avec variable X
\usepackage{awesomebox}%---Pour les boites info, danger et autres
\usepackage{menukeys}%---Pour deiffner les touches de Calculatrice
\usepackage{fancyhdr}%---pour les en-t\^ete personnalis\'ees
\usepackage{blindtext}%---pour les liens
\usepackage{hyperref}%---pour les liens (\`a mettre en dernier)
\usepackage{caption}%---pour la francisation de la l\'egende table vers Tableau
\usepackage{pifont}
\usepackage{array}%---pour les tableaux
\usepackage{lipsum}
\usepackage{yFlatTable}
\usepackage{multicol}
\newcommand{\Lim}[1]{\lim\limits_{\substack{#1}}\:}
\renewcommand{\vec}{\overrightarrow}
\newcommand{\N}[0]{\mathbb{N}}
\newcommand{\dd}{\mathrm{d}}
\newcommand{\norm}[1]{||\vec{#1}||}
\newcommand{\fo}{\psi(\vec{r},t)}
\newcommand{\foe}{\psi(\vec{r},t)\*}
\newcommand{\HH}{\hat{H}}
\newcommand{\hb}{\hbar}
\newcommand{\lap}{\nabla^2}
\newcommand{\lapcc}{\frac{\partial^2 }{\partial x^2}+\frac{\partial^2 }{\partial y^2}+\frac{\partial^2 }{\partial z^2}}
\newcommand{\E}{\vec{E}(M)}
\newcommand{\B}{\vec{B}(M)}
\newcommand{\V}{V(M)}
\newcommand{\er}{\vec{e_{\rho}}}
\newcommand{\ep}{\vec{e_{\varphi}}}
\newcommand{\pip}{\pi^+}
\newcommand{\pim}{\pi^-}
\newcommand{\mv}{\mathcal{V}}
\begin{document}
\setcounter{chapter}{-1}
\chapter{Électromagnétisme : Formulaire}
\section{Champ électrique}
\subsection{Charge ponctuelle}
\begin{definition}[Champ électrique]
\[\vec{F_{q_p\to q}} = q \cdot \E\] avec \(\E\) qui ne dépend que de la charge \(q_p\) : \[\E = \frac{q_p}{4\pi\epsilon_0}\frac{\vec{e_{PM}}}{\norm{PM}^2}\] est le champ électrique crée par \(q_p\) au point M, en \(NC^{-1} / Vm^{-1}\)
\end{definition}
\begin{definition}[Potentiel électrique]
\(V(M) = \frac{q}{4\pi\epsilon_0 r} + Cst\)
\end{definition}
\subsection{Équation entre E et V entre E et V}
\begin{theorem}[Équation intégrale]
\[V(B)-V(A) = -\int_A^B\E\cdot \dd \vec{l}(M)\]
\end{theorem}
\begin{theorem}[Équation locale entre E et V]
\[\dd V = -\E\cdot \dd \vec{l}(M)\]
\[\E = -grad(\V)\]
\end{theorem}
\subsection{Distribution volumique de charge}
\begin{theorem}[Champ électrique total  (d. volumique)]
\[\E = \iiint_{P\in\mv} \frac{\rho_c(P)\dd \mv(P)\vec{PM}}{4\pi\epsilon_0 \norm{PM}^3}\]
 \end{theorem}
 \begin{theorem}[Potentiel total  (d. volumique)]
\[V(M) = \iiint_{P\in \mv} \frac{\rho_c \dd \mv(P)}{4\pi\epsilon_0 \norm{PM}}\]
 \end{theorem}
\subsection{Théorème de Gauss}
\begin{theorem}[Théorème de Gauss]
Le flux de E à travers une surface fermée orientée est égal à la somme des charges intérieures à cette surface fermée divisée par \(\epsilon_0\) :
\[\iint \E\cdot \vec{n}\dd S = \frac{\sum q_{int}}{\epsilon_0}\]
\end{theorem}
\subsection{Symétrie}
\begin{center}
\begin{tabular}{ll}
Situation & Direction\\
M \(\in \pip\) & E est contenu au plan\\
M \(\in \pim\) & B est \(\perp\) au plan\\
\( (Oxy) \in  \pip\) & Selon z, \(E_z\) impaire, \(Bx,B_y\) paires\\
\( (Oxy) \in  \pim\) & Selon z, \(E_z\) paire, \(Bx,B_y\) impaires
\end{tabular}
\end{center}
\section{Champ magnétique}
\subsection{Symétrie}
\begin{center}
\begin{tabular}{ll}
Situation & Direction\\
M \(\in \pip\) & B est \(\perp\) au plan\\
M \(\in \pim\) & B est contenu dans le plan\\
\( (Oxy) \in  \pip\) & Selon z, \(B_z\) paire, \(Bx,B_y\) impaires\\
\( (Oxy) \in  \pim\) & Selon z, \(B_z\) impaire, \(Bx,B_y\) paires
\end{tabular}
\end{center}
\subsection{Théorème d'Ampère}
\begin{theorem}[Théorème d'Ampère]
\[\int_C \vec{B}(M)\cdot \dd \vec{l} = \mu_0 \sum I_{\text{entrées par C}} \]
Les courant sont comptés positivement s'ils sont dans le sens de \(\vec{n}\).
\end{theorem}

\begin{theorem}[Théorème d'Ampère pour une distribution volumique de courant]
 \[\int_C\B\cdot \dd \vec{l} = \mu_0 \iint \vec{j}\cdot \vec{n}\dd S\]
\end{theorem}

\begin{theorem}[Forme locale du théorème d'Ampère]
\[Rot(\vec{B})=\mu_0 \vec{j}\]
\end{theorem}
\subsection{Loi de Biot et Savat}
\begin{definition}[Intensité du courant électrique]
\[I = \iint_{S}\vec{j}\cdot \vec{n}\dd S\] avec S la section étudiée
\end{definition}
\begin{theorem}[Loi de Biot et Savart]
 \[\B = \frac{\mu_0}{4\pi}  \int_{P\in F} \frac{I\vec{\dd l}(P)\wedge \vec{PM}}{\norm{PM}^3}\]
\end{theorem}
\subsection{Résultats à savoir}
\begin{proposition}[Champ crée par un fil infiniment long]
\[\vec{B}(M) = \frac{\mu_0 I}{2\pi \rho}\vec{e_{\varphi}}\]
\end{proposition}
\begin{proposition}[Champ crée par une spire]
\[\B = \frac{\mu_0 I}{2R}\sin^3(\alpha)\vec{e_z}\] avec \(\alpha\) l'angle sous lequel on voit la spire depuis M.
\end{proposition}
 \end{document}
